\documentclass[11pt,a4paper]{article}
\usepackage[T1]{fontenc}
\usepackage{lmodern}
\usepackage[margin=24mm,headheight=15pt]{geometry}
\usepackage{amsmath,amssymb,amsthm,mathtools,booktabs,array,tabularx}
\usepackage{microtype,fancyhdr,xcolor}
\usepackage[unicode]{hyperref}
\definecolor{ink}{RGB}{28,61,79}
\hypersetup{colorlinks=true,linkcolor=ink,urlcolor=ink,citecolor=ink,
 pdftitle={SRG Signed-Grid Reachability: Exact Two-Step Access and Sharp Small-Grid Thresholds},
 pdfauthor={Hilmir Frímann Halldórsson}}
\pagestyle{fancy}\fancyhf{}
\fancyhead[L]{\small SRG Signed-Grid Reachability}
\fancyhead[R]{\small v0.1 | Manuscript for review}
\fancyfoot[C]{\thepage}
\setlength{\parindent}{0pt}\setlength{\parskip}{5pt plus 1pt minus 1pt}
\setlength{\emergencystretch}{2em}
\clubpenalty=10000 \widowpenalty=10000
\allowdisplaybreaks[1]
\newcolumntype{Y}{>{\raggedright\arraybackslash}X}
\theoremstyle{definition}
\newtheorem{theorem}{Theorem}[section]
\newtheorem{proposition}[theorem]{Proposition}
\newtheorem{lemma}[theorem]{Lemma}
\theoremstyle{remark}\newtheorem{remark}[theorem]{Remark}
\DeclareMathOperator{\sgn}{sgn}
\newcommand{\R}{\mathbb R}
\numberwithin{equation}{section}
\title{\textbf{SRG Signed-Grid Reachability}\\[4pt]
\large Exact Two-Step Access and Sharp Small-Grid Thresholds}
\author{Hilmir Frímann Halldórsson}
\date{Version v0.1 --- 12 October 2026\\[3pt]\small Manuscript and proposed publication payload awaiting review}
\begin{document}
\maketitle
\begin{abstract}
For January's finite periodic field, signed compression leaves a gap between zero
and the smallest nonzero output magnitude. We reproduce the accepted analysis of
how this gap constrains access to signed floor targets at the second pass.
The uncompressed map has an ordinary global smooth inverse, while the compressed
map has set-valued preimages. Their distinction yields an exact intersection
criterion and a complete precursor formula for specified positive gains.
On a singleton the second-gain threshold is approximately $0.889066849332978844$.
For the source's actual anisotropic $1\times2$ stencil, sharp same-sign and
opposite-sign thresholds are approximately $1.189427271430029358$ and
$1.796277920195400199$. Complete sign-sector proofs include unequal magnitudes,
ceiling values and boundary cases. At unit gain, a zero-valued site can receive its
neighbor's contribution and reach a same-sign floor target even though that target
does not remain fixed. These are exact-real existential access results, not
claims about source-schedule realization, probabilities or later dynamics.
\end{abstract}

\section{A zero site can receive its neighbor}\label{sec:opening}
Consider two sites on the periodic second axis of the January grid. Write
$u=e^{-10}$ for its positive compression floor, and temporarily take the second
gain to be one. A zero at one site does not prevent a contribution from the other:
the source stencil has a positive off-diagonal coefficient $5/18$.

\begin{center}
\begin{tabular}{@{}lll@{}}
\toprule Intermediate field & Before compression & After signed compression\\
\midrule
$(u,0)$ & $(A,B)$, with $0<B<A<u-10^{-6}$ & $(u,u)$\\
$(u,u)$ & $(u+\tfrac1{10}\tanh u)(1,1)$ & strictly above $(u,u)$\\
$(u,-u)$ & $(D,-D)$, with $0<D<u-10^{-6}$ & $(u,-u)$\\
\bottomrule
\end{tabular}
\end{center}
Here $A=(13/18)u+\tanh(u)/10$, $B=(5/18)u$ and
$D=(4/9)u+\tanh(u)/10$. We prove every inequality in the table below.
The first row demonstrates access to a target; the second shows that reaching
the target and remaining there are different questions. Neither row is a simulated
trajectory. It specifies the field entering the second pass.

To call this \emph{two-step} access, one must also ask whether the intermediate
field can emerge from a first pass. This is where the signed compressor's complete
range and the uncompressed map's inverse matter. The first pass can reach every
field whose coordinates lie in its scalar range, but cannot put a nonzero coordinate
strictly between $-u$ and $u$. That gap yields a genuine obstruction on a singleton.
Coupling changes the sharp obstruction on the actual $1\times2$ grid.

\section{The January finite periodic field}\label{sec:source}
We use the historical January operators, at the specified retained parameters,
in their exact-real interpretation. This is not Tower II's moving-particle model.
The frozen source files J1--J3 \cite{J1,J2,J3}, the accepted OC02 inverse framework
\cite{G}, the source/notation addendum \cite{A}, and the accepted OC06 and OC08
reports \cite{OC06,OC08} determine the account.
Line references appear in Appendix~\ref{app:sources}; full identities and access
limitations are in the companion source index.

For an $N\times M$ real array $X$, with indices taken modulo $N,M$, the original
anisotropic periodic-index operator at $a=1$, $b=3/5$ is
\begin{equation}\label{eq:L}
\begin{split}
 (LX)_{ij}={}&X_{i+1,j}+X_{i-1,j}-2X_{ij}\\
 &+\frac{25}{9}\bigl(X_{i,j+1}+X_{i,j-1}-2X_{ij}\bigr).
\end{split}
\end{equation}
Repeated neighbors are repeated \emph{incidences}. For an axis of length one its
two self-incidences cancel its $-2X$ term; for length two the same other site is
counted twice. The first axis is the source's $\theta$ axis and the second its
$\varphi$ axis. The displayed embedding does not turn this index stencil into a
surface Laplace--Beltrami operator. In particular, transposing $1\times2$ into
$2\times1$ would change the coefficients used here.

With $h=1/20$, nonlinear coefficient $1/10$ and scale one, flatten the array to
$\R^d$, $d=NM$, and set
\begin{equation}\label{eq:P}
 T=I+\frac1{20}L,\qquad
 P(X)=TX+\frac1{10}\tanh(X),\qquad F_m(X)=C_s(mP(X)),\quad m>0.
\end{equation}
The hyperbolic tangent is coordinatewise and uses the incoming field. It is
\emph{not} multiplied by another factor $h$.
The gain $m$ is specified as a positive parameter. No claim is made that arbitrary
successive values $m_1,m_2$ occur in the historical schedule.

This analysis explicitly selects signed compression. The unchanged application's
initial menu selection is paper mode, not signed mode. The app offers sizes
$64,128,256$; the small arrays studied here are mathematically admissible under
the periodic functions, not claims about those menu choices.
The relevant app defaults are $L_c=10$, $\epsilon=10^{-6}$; the generic compressor
function has different defaults, which are not substituted into this paper.
Define
\begin{equation}\label{eq:Cs}
\begin{gathered}
 u=e^{-10},\qquad v=e^{10},\qquad f=u-\epsilon>0,\\
 C_s(q)=
 \begin{cases}
 0,&q=0,\\
 \sgn(q)\min\{v,\max\{u,|q|+\epsilon\}\},&q\ne0.
 \end{cases}
\end{gathered}
\end{equation}
The positivity of $f$ follows, for example, from $e<3$:
$u>1/59049>10^{-6}$. Exact zero is a separate branch.
Statements about real inputs and exact endpoint inclusions do not imply that all
these real values are exactly representable by the historical floating-point code.

\section{The ordinary inverse belongs to P}\label{sec:inverse}
The following is the accepted OC02 inverse theorem, reproduced with the reasoning
needed here. It is neither a new theorem-suite execution nor a novelty claim.

\begin{theorem}[Global smooth inverse of the retained $P$]\label{thm:inverse}
For every finite positive pair of grid dimensions $N,M$, the map $P$ in
\eqref{eq:P} is a smooth bijection $\R^d\to\R^d$ with smooth inverse.
\end{theorem}
\begin{proof}
Periodic shifts and their inverses occur in pairs, so $L$ and $T$ are symmetric.
The discrete Fourier modes have $T$ eigenvalues
\begin{equation}\label{eq:eigen}
 \lambda_{pq}=1-\frac15\sin^2\!\frac{\pi p}{N}
                 -\frac59\sin^2\!\frac{\pi q}{M}
 \ \geq\ 1-\frac15-\frac59=\frac{11}{45}>0.
\end{equation}
This formula includes lengths one and two, with their actual repeated incidences.
Thus $T$ is positive definite. Furthermore,
\begin{equation}\label{eq:DP}
 DP(X)=T+\frac1{10}\operatorname{diag}(\operatorname{sech}^2 X_i)
 \succeq\frac{11}{45}I.
\end{equation}
Write
\begin{equation}\label{eq:potential}
 \Psi(X)=\frac12 X^{\top}TX+\frac1{10}\sum_{i=1}^d\log\cosh X_i .
\end{equation}
Then $\nabla\Psi=P$. For each $z\in\R^d$, the function
$\Psi(X)-z^{\top}X$ is strictly convex and coercive, because
$\log\cosh X_i\geq0$ and its quadratic term bounds it below by
$(11/90)\|X\|^2-\|z\|\|X\|$. It has a unique minimizer, whose gradient equation
is $P(X)=z$. This proves both surjectivity and injectivity.
The Jacobian is invertible everywhere by \eqref{eq:DP}. The inverse function
theorem supplies local smooth inverses, which agree as the single global inverse
because of bijectivity.
\end{proof}

Compression changes this conclusion. Entire intervals map to a floor or ceiling,
so $F_m$ is not injective. We write $P^{-1}$ for an ordinary inverse and use
explicit sets for compressed preimages. Confusing them would lose entire families
of precursors.

\section{The complete signed range and scalar fibers}\label{sec:range}
Let
\begin{equation}\label{eq:S}
 S=\{0\}\cup[-v,-u]\cup[u,v],\qquad J(y)=\{q\in\R:C_s(q)=y\}.
\end{equation}

\begin{proposition}[Scalar range and all preimages]\label{prop:fibers}
The image of $C_s$ is exactly $S$. Its full scalar fibers are:
\begin{center}
\begin{tabular}{@{}ll@{}}
\toprule Output $y$ & $J(y)$\\
\midrule
$0$ & $\{0\}$\\
$u$ & $(0,f]$\\
$-u$ & $[-f,0)$\\
$u<|y|<v$ & $\{\,y-\epsilon\sgn(y)\,\}$\\
$v$ & $[v-\epsilon,\infty)$\\
$-v$ & $(-\infty,-(v-\epsilon)]$\\
$y\notin S$ & $\varnothing$\\
\bottomrule
\end{tabular}
\end{center}
\end{proposition}
\begin{proof}
Every nonzero input keeps its sign and has compressed magnitude in $[u,v]$,
while zero maps to zero. This proves the forward range inclusion.
Conversely, for $u\leq|y|\leq v$, choose
$q=\sgn(y)(|y|-\epsilon)$. Since $u>\epsilon$, this input is nonzero, has the
correct sign, and satisfies $|q|+\epsilon=|y|$; hence $C_s(q)=y$.
Together with $q=0$ this proves surjectivity onto $S$, including all endpoints.

For positive nonzero input, output $u$ occurs exactly when
$q+\epsilon\leq u$, or $0<q\leq f$. A strictly interior output requires
$q+\epsilon=y$ uniquely, and output $v$ requires $q+\epsilon\geq v$.
The analogous negative cases follow from oddness. No nonzero input produces
zero, and the range proof excludes all remaining outputs. These observations
give each row of the table with its stated endpoint inclusion.
\end{proof}

\begin{proposition}[First-pass range]\label{prop:range}
For every $m>0$,
\begin{equation}\label{eq:range}
 F_m(\R^d)=S^d.
\end{equation}
\end{proposition}
\begin{proof}
The scalar range gives $F_m(\R^d)\subseteq S^d$. If $Y\in S^d$, choose any
$z\in\prod_i J(Y_i)$; every factor is nonempty by Proposition~\ref{prop:fibers}.
Theorem~\ref{thm:inverse} supplies $X=P^{-1}(z/m)$, for which
$C_s(mP(X))=C_s(z)=Y$. This proves the reverse inclusion.
\end{proof}
This is an existence statement over unrestricted real incoming fields. It neither
samples a precursor nor proves its probability under an initializer.

\section{An exact two-step criterion and every precursor}\label{sec:twostep}
Fix positive gains $m_1,m_2$ and a sign field $\sigma\in\{-1,1\}^d$.
The target is the signed floor vector $\sigma u$.
Define its second-pass capture cell
\begin{equation}\label{eq:B}
 B_{m_2,\sigma}
 =\{Y\in\R^d:0<\sigma_i m_2P_i(Y)\leq f\text{ for every }i\}.
\end{equation}
The lower face is excluded: exact zero is not a signed floor.
The upper face is included: a magnitude of $f$ becomes $u$ after adding
$\epsilon$.

\begin{theorem}[Two-step access and full precursor set]\label{thm:twostep}
There exists $X$ with $F_{m_2}(F_{m_1}(X))=\sigma u$ if and only if
\begin{equation}\label{eq:intersection}
 S^d\cap B_{m_2,\sigma}\ne\varnothing.
\end{equation}
The complete set of such initial fields is
\begin{equation}\label{eq:all}
\begin{split}
 \mathcal A_{m_1,m_2,\sigma}
 =\bigcup_{Y\in S^d\cap B_{m_2,\sigma}}
 \left\{P^{-1}(z/m_1):
             z\in\prod_{i=1}^d J(Y_i)\right\}.
\end{split}
\end{equation}
\end{theorem}
\begin{proof}
If $X$ succeeds, set $Y=F_{m_1}(X)$. Proposition~\ref{prop:range} gives $Y\in S^d$,
and the signed floor fibers give $Y\in B_{m_2,\sigma}$.
Moreover $z=m_1P(X)$ lies in $\prod_iJ(Y_i)$, so $X$ belongs to the stated union.
Conversely, any member of the union has first image $Y$, by its scalar fiber
conditions, and second image $\sigma u$, by \eqref{eq:B}. If the intersection
is nonempty, its scalar fibers are nonempty and the ordinary inverse $P^{-1}$
constructs a member. This proves both the equivalence and completeness.
\end{proof}

Nonemptiness is independent of the first positive gain. The actual precursor
fields and their geometry can depend on $m_1$, as \eqref{eq:all} shows.
The theorem concerns being at the target at pass two, not necessarily arriving
there for the first time. It asserts neither that the target is fixed nor that
the intersection has any specified volume or initialization probability.

For comparison, one-pass access to any signed floor target exists from an unrestricted
real input: choose $z_i=\sigma_i f/2$ and $X=P^{-1}(z/m)$.
The nontrivial second-pass obstruction arises because an intermediate must also
belong to $S^d$. That restriction cannot be dropped by invoking surjectivity of $P$.

\section{A singleton: a sharp gap obstruction}\label{sec:singleton}
For a $1\times1$ array, both periodic self-incidences cancel on both axes.
Thus $L=0$ and
\begin{equation}\label{eq:singleP}
 P(y)=y+\frac1{10}\tanh y.
\end{equation}
This is odd and strictly increasing, with derivative
$1+\operatorname{sech}^2(y)/10>0$. For positive $y$ it exceeds $y$.

\begin{theorem}[Singleton threshold]\label{thm:singleton}
For every specified $m_1>0$, either singleton signed floor target is accessible at
pass two exactly when
\begin{equation}\label{eq:singlethreshold}
 0<m_2\leq m_{\rm single}
 :=\frac{f}{u+\tanh(u)/10}.
\end{equation}
\end{theorem}
\begin{proof}
An intermediate in $S$ is either zero or has magnitude at least $u$.
Zero remains zero and cannot reach a signed target. For any other intermediate,
oddness and strict increase imply $|P(y)|\geq P(u)>0$.
The signed floor condition therefore requires $m_2P(u)\leq f$.
Conversely, if this inequality holds, use $y=u$ for the positive target or
$y=-u$ for the negative target. The first-pass range reaches that $y$ for
every $m_1>0$. At equality the second precompression magnitude is exactly $f$,
which is included in the floor fiber.
\end{proof}

At gain one access is impossible, since $P(u)>u>f$.
Gain $1/2$ does permit access: $\tanh u<u$ gives
$0<P(u)/2<11u/20$, while
\begin{equation}\label{eq:halfbound}
 f-\frac{11u}{20}
 =\frac9{20}u-\epsilon
 >\frac9{20\cdot59049}-10^{-6}
 =\frac{43439}{6561000000}>0.
\end{equation}
Here $e<3$ supplies $u>1/59049$; no decimal evaluation is needed.
The saved accepted OC06 threshold, rounded to 18 places, is
$m_{\rm single}\approx0.889066849332978844$.
No initial-field inverse was numerically solved to obtain a witness.

\section{The actual coupled $1\times2$ reduction}\label{sec:twosite}
Let the first axis have length one and the second length two.
The first-axis stencil cancels. On the second axis, forward and backward neighbors
coincide but count twice. Consequently
\begin{equation}\label{eq:Ttwo}
 L=\frac{50}{9}
 \begin{pmatrix}-1&1\\1&-1\end{pmatrix},
 \qquad
 T=\begin{pmatrix}13/18&5/18\\5/18&13/18\end{pmatrix}.
\end{equation}
Define, only for this January matrix,
\begin{equation}\label{eq:alphabeta}
 \alpha=\frac{13}{18},\quad\beta=\frac5{18},\quad
 \delta=\alpha-\beta=\frac49,\quad h(y)=\frac1{10}\tanh y.
\end{equation}
Then $\alpha+\beta=1$, and
\begin{equation}\label{eq:Ptwo}
 P_1(Y)=\alpha Y_1+\beta Y_2+h(Y_1),\qquad
 P_2(Y)=\beta Y_1+\alpha Y_2+h(Y_2).
\end{equation}
The symbol $h(y)$ here names the scalar nonlinearity; the fixed time-step coefficient
in \eqref{eq:P} remains $1/20$.

\begin{lemma}[Two symmetries]\label{lem:symmetry}
Global sign reversal and exchange of the two sites preserve $S^2$ and commute
with $P$, coordinatewise $C_s$, and therefore $F_m$.
\end{lemma}
\begin{proof}
The set $S$ is sign symmetric, so global negation preserves $S^2$;
coordinate exchange also preserves a Cartesian square.
The nonlinearity is odd and the linear map changes sign under global negation,
so $P(-Y)=-P(Y)$. Equation~\eqref{eq:Ptwo} exchanges its two components when
$Y_1,Y_2$ are exchanged. The same operations commute with a common odd scalar
compressor applied to both coordinates. This proves each assertion.
\end{proof}
It is thus enough to prove the same-sign case $(u,u)$ and the opposite-sign
case $(u,-u)$. The other two targets follow analytically, not from additional
sampled fixtures.

\section{Same-sign access: exhaustive necessity}\label{sec:same}
Put
\begin{equation}\label{eq:AB}
 A=\alpha u+h(u),\qquad B=\beta u.
\end{equation}
Both are positive and $B<A$, since
$A-B=\delta u+h(u)>0$.

\begin{theorem}[Sharp same-sign threshold]\label{thm:same}
For every $m_1>0$, a same-sign two-site floor target is accessible at pass two
if and only if
\begin{equation}\label{eq:msame}
 0<m_2\leq m_{\rm same}:=\frac fA.
\end{equation}
\end{theorem}
\begin{proof}
By Lemma~\ref{lem:symmetry} it suffices to consider the positive target.
Any admissible intermediate $Y\in S^2$ must have both $P_i(Y)>0$ and
$\max_iP_i(Y)\leq f/m_2$. We exhaust all sectors.

If both coordinates are nonpositive, the nonnegative matrix coefficients and
odd increasing $h$ make both $P_i$ nonpositive, so this sector is excluded.
If exactly one coordinate is positive and the other zero, write it, up to
exchange, as $(a,0)$ with $u\leq a\leq v$. Then
$P(Y)=(\alpha a+h(a),\beta a)$; the first is the larger and is at least $A$.
The smallest value is attained at $a=u$.
If both coordinates are positive, each is at least $u$, so each output is at least
$(\alpha+\beta)u+h(u)=u+h(u)>A$.

It remains to examine opposite signs without assuming equal magnitudes.
By exchange write $Y=(a,-b)$ with $u\leq a,b\leq v$.
If $P_2\leq0$, this field cannot reach the positive target.
If $P_2>0$, eliminating $a$ from \eqref{eq:Ptwo} gives the exact identity
\begin{equation}\label{eq:eliminate}
 P_1=\frac{\alpha}{\beta}P_2+
       \frac{\alpha^2-\beta^2}{\beta}b+
       h(a)+\frac{\alpha}{\beta}h(b).
\end{equation}
All terms except the displayed multiple of $b$ are strictly positive, and
$(\alpha^2-\beta^2)/\beta=8/5$. Hence
$P_1>(8/5)u>A$: indeed $\tanh u<u$ implies
$A<(13/18+1/10)u=(37/45)u<(8/5)u$.
This includes unequal magnitudes and the ceiling $v$. All sectors have now
been covered. Any successful intermediate therefore satisfies $\max_iP_i\geq A$,
which forces $m_2A\leq f$.

Conversely take $Y=(u,0)\in S^2$. Its outputs are $(A,B)$ with $0<B<A$.
If $m_2A\leq f$, both lie in the positive capture interval after multiplication
by $m_2$. At $m_2=f/A$, the first equals $f$ and the second equals $fB/A\in(0,f)$.
The first-stage range theorem supplies this intermediate for every $m_1>0$.
The negative target follows by global sign reversal.
\end{proof}

The zero in this optimal witness is not an absence of coupling. It is exactly
the positive term $\beta u$ entering that site that makes the target accessible.
The proof establishes a sharp bound over \emph{all} intermediates in $S^2$,
not merely over a chosen family.

\section{Opposite-sign access without an antisymmetry assumption}\label{sec:opposite}
Define
\begin{equation}\label{eq:D}
 D=\delta u+h(u)>0 .
\end{equation}
\begin{theorem}[Sharp opposite-sign threshold]\label{thm:opposite}
For every $m_1>0$, an opposite-sign floor target is accessible at pass two
if and only if
\begin{equation}\label{eq:mopp}
 0<m_2\leq m_{\rm opposite}:=\frac fD.
\end{equation}
\end{theorem}
\begin{proof}
Consider $(u,-u)$; its required signed outputs are $P_1>0$ and $-P_2>0$.
If both coordinates are zero, both outputs are zero.
With just one nonzero coordinate, the positive off-diagonal coefficient and the
same-sign diagonal/nonlinear terms make both outputs have that coordinate's sign.
With both coordinates of the same sign, both outputs again have that sign.
These sectors cannot give the required opposite signs.

A reversed-sign intermediate has $Y=(-a,b)$, $a,b\in[u,v]$. It obeys
$P_1-P_2=-\delta(a+b)-h(a)-h(b)<0$, contradicting $P_1>0>P_2$.
The only remaining possibility is the aligned sector $Y=(a,-b)$.
For it set
\[
 p_1=P_1=\alpha a-\beta b+h(a),\qquad
 p_2=-P_2=\alpha b-\beta a+h(b).
\]
If either is nonpositive, capture fails. Otherwise,
\begin{equation}\label{eq:average}
 \max(p_1,p_2)\geq\frac{p_1+p_2}{2}
 =\frac{\delta(a+b)+h(a)+h(b)}2
 \geq\delta u+h(u)=D.
\end{equation}
The last inequality uses $a,b\geq u$ and monotonicity of $h$.
It makes no assumption that $a=b$, and includes the ceiling endpoints.
Necessity of $m_2D\leq f$ follows.

For sufficiency use $Y=(u,-u)\in S^2$, for which $P(Y)=(D,-D)$.
If $m_2D\leq f$, both signed capture inequalities hold.
At the upper endpoint both required magnitudes equal $f$, which is allowed.
The first-stage range theorem supplies this intermediate for every $m_1>0$.
Global reversal or site exchange yields the other target.
\end{proof}
Antisymmetry supplies a witness attaining the bound; it was not imposed on the
necessity argument. This distinction is essential to sharpness over the complete
admissible range.

\section{Reaching, remaining and complete witness precursors}\label{sec:interpretation}
\subsection{Unit gain and the exact margin}
The inequality $\tanh u<u$ gives $A<(37/45)u$, and therefore
\begin{equation}\label{eq:unitmargin}
 f-A>\frac8{45}u-\epsilon
 >\frac8{45\cdot59049}-10^{-6}
 =\frac{1068559}{531441000000}>0.
\end{equation}
Thus $0<B<A<f$, proving $F_1(u,0)=(u,u)$.
Since $0<D<A$, the same bound proves $F_1(u,-u)=(u,-u)$.
But $P(u,u)=(u+h(u))(1,1)$ is strictly above $u$ in each coordinate.
After adding $\epsilon$ the compressed magnitude is still above $u$:
if clipping at the ceiling occurred it would be $v>u$, so either branch suffices.
Consequently $(u,u)$ is not fixed at gain one.
The positive value received at the initially zero site is $B=\beta u$,
not a compression convention that changes zero directly to $u$.

\subsection{All first-stage precursors of the witnesses}
For arbitrary specified $m_1>0$, all initial fields whose first image is $(u,0)$ are
\begin{equation}\label{eq:precsame}
 \left\{P^{-1}\bigl((z/m_1,0)\bigr):0<z\leq f\right\}.
\end{equation}
This is the complete preimage of that intermediate by the scalar fiber table.
The zero is a coordinate of the intermediate and of its precompression vector.
Applying the coupled inverse $P^{-1}$ does not imply that the corresponding
coordinate of the initial field is zero.
All initial fields whose first image is $(u,-u)$ are
\begin{equation}\label{eq:precopp}
 \left\{P^{-1}\bigl((z_1/m_1,z_2/m_1)\bigr):
               0<z_1\leq f,\ -f\leq z_2<0\right\}.
\end{equation}
These families give successful two-step precursors whenever the respective
second-gain threshold holds. They need not exhaust the other successful intermediate
fields. The exhaustive description for the target is the union \eqref{eq:all}.
No inverse equation was solved numerically and no basin volume follows from
writing these sets.

\subsection{What the three thresholds compare}
Exact denominator inequalities give
\begin{equation}\label{eq:thresholdorder}
 0<D<A<u+h(u),\qquad
 m_{\rm single}<m_{\rm same}<m_{\rm opposite},\qquad
 m_{\rm single}<1<m_{\rm same}.
\end{equation}
Here $A-D=\beta u>0$, $u+h(u)-A=(1-\alpha)u>0$;
$m_{\rm single}<1$ follows from $f<u<u+h(u)$, and $m_{\rm same}>1$
from \eqref{eq:unitmargin}.

\begin{center}
\begin{tabular}{@{}llr@{}}
\toprule Grid and target & Exact upper second gain & Saved value, 18 places\\
\midrule
$1\times1$, either sign & $f/[u+h(u)]$ & $0.889066849332978844$\\
$1\times2$, same signs & $f/[\alpha u+h(u)]$ & $1.189427271430029358$\\
$1\times2$, opposite signs & $f/[\delta u+h(u)]$ & $1.796277920195400199$\\
\bottomrule
\end{tabular}
\end{center}
These are thresholds for the stated grids and target classes, not for arbitrary
arrays. In particular, the opposite-sign number exceeding $3/2$ is a result about
specified positive gains, not about values attained by the source's historical
gain schedule.

\section{Evidence and limits}\label{sec:evidence}
The central distinction is between unrestricted one-pass capture, two-pass access
through $S^d$, and fixedness of a target. Later persistence, first arrival,
initialization probability and finite-precision behavior are further questions
not resolved by these theorems. No full-grid classification, long-time dynamics,
new recursion, source-schedule matching, entropy, cosmology or physical validation
is established. No Tower amplitude theorem is used as a January premise.
This manuscript is not the reserved SRG-II/R0--R3 mechanism paper.

The mathematical inputs OC06, OC08 and the inherited OC02 framework are accepted
within their declared scopes. The new v0.1 manuscript and six-file candidate
await manuscript-and-payload review. The human author is Hilmir Frímann Halldórsson.
GPT supplied work orders and scientific review of the accepted analyses;
Codex prepared this manuscript, checked static source/proof correspondence,
transcribed saved decimals, compiled it and inspected the pages. AI assistance
is not external human peer review.

The historical final OC06 record contains eight passing check groups and OC08
contains seven. No accepted checker or supplied GPT review script was rerun,
and no new scalar calculation, inverse solve, grid enumeration, simulation or
source execution was performed for this manuscript.
The three numerical thresholds in the abstract and comparison table were
transcribed from saved JSON strings and rounded to 18 decimal places by decimal
round-half-even. The exact inequalities, not those decimals, prove the conclusions.
The companion evidence summary and private review records identify the saved fields.

The source index distinguishes original public source routes from retained private
reports. Complete definitions and proofs are given here so that the argument
does not require access to private chat or reports. No new blanket reuse licence
is assigned to this manuscript or reconstruction records. Existing source-specific
CC BY 4.0 notices remain attached to their originals; the software licence is
not presumed to cover this research.

\appendix
\section{Native source correspondence}\label{app:sources}
Line numbers below are one-based in the frozen complete source texts.
The byte-order mark in J2 is included in its complete-file identity, even though
it is removed when decoding for reading. The private review bundle includes
only a narrow derivative excerpt of the consumed lines; it is neither a complete
source nor a substitute for whole-file verification.

\begin{center}\small
\begin{tabularx}{\textwidth}{@{}llY@{}}
\toprule Source & Lines & Definition used\\
\midrule
J1 & 4--5, 19--40 & periodic rolls; axes; $1/a^2,1/b^2$ stencil weights\\
J1 & 62--76 & display embedding and axis interpretation\\
J2 & 37--67 & generic compressor defaults and signed branch\\
J2 & 109--138 & incoming field, TGMO step, separate nonlinearity, gain, compression\\
J3 & 44--66 & offered sizes and app defaults, including initial paper mode\\
J3 & 234--238 & incoming-field nonlinear callback\\
G & 238--252 & accepted global inverse proof for $P$\\
A & 13--31 & January notation and source distinctions\\
\bottomrule
\end{tabularx}
\end{center}
The exact source equations, rather than an assumed embedded-surface geometry,
fix the distinction between the $1\times2$ and $2\times1$ reductions.
OC06 supplies the range, fiber, two-step and singleton arguments; OC08 supplies
the complete two-site sign-sector proofs and attaining witnesses.

\begin{thebibliography}{9}
\bibitem{J1} Historical January source, \emph{tgmo.py}, frozen J1, 3,208 bytes.
Public source route: \url{https://zenodo.org/records/18285147}, file
\texttt{tgmo.py}. Exact SHA-256 in the companion source index.
\bibitem{J2} Historical January source, \emph{refu.py}, frozen J2, 6,495 bytes
including UTF-8 BOM. Same public record, file \texttt{refu.py}.
\bibitem{J3} Historical January source, \emph{app.py}, frozen J3, 14,800 bytes.
Same public record, file \texttt{app.py}. This is the retained version, not a
substitution from record 18366687.
\bibitem{G} \emph{OC02 Signed Plateaus}, accepted report G, 11 October 2026.
Retained final handoff, \path{updates/OC02_SIGNED_PLATEAUS_REPORT.md}.
\bibitem{A} \emph{Required Source and Notation Addendum}, retained final handoff,
\path{updates/REQUIRED_SOURCE_AND_NOTATION_ADDENDUM.md}.
\bibitem{OC06} \emph{OC06 January Two-Step Access}, accepted REPORT.md and
CHECK\_RESULTS.json, 12 October 2026. Retained private research records.
\bibitem{OC08} \emph{OC08 January Two-Site Access}, accepted REPORT.md and
CHECK\_RESULTS.json, 12 October 2026. Retained private research records.
\end{thebibliography}
\end{document}
